\documentclass[11pt]{article}
\usepackage[margin=1in]{geometry}
\usepackage{graphicx}
\usepackage{booktabs}
\usepackage{amsmath,amssymb}
\usepackage{caption}
\usepackage{float}
\usepackage{xcolor}
\usepackage[hidelinks]{hyperref}
\usepackage{etoolbox}
% Start every top-level section on a new page
\pretocmd{\section}{\clearpage}{}{}

\newcommand{\code}[1]{\texttt{#1}}
\newcommand{\kt}{\kappa_t}
\newcommand{\Nbar}{\bar N}
\newcommand{\Nhat}{\hat N}
\graphicspath{{/Users/satoshi/GitHub/NKPC_HSA_MCMC/results/figures/}}

\title{\textbf{Market Concentration and the Flattening of the Phillips Curve}\\[2pt]
\large A Theory-Motivated, Semi-Structural State-Space Assessment of the HSA Mechanism}
\author{Satoshi Ichikawa}
\date{Draft --- July 2026}

\begin{document}
\maketitle

\noindent\fcolorbox{red}{red!5}{\begin{minipage}{0.94\linewidth}
\textbf{Re-estimation required.} The saved numerical results below predate the
state-initialization and restricted-covariance corrections. They are not used for
conclusions until rerun with
\code{estimation\_revision=2026-08-state-initial-covariance-v2}.
\end{minipage}}

\begin{abstract}
\noindent We ask whether the slope of the New Keynesian Phillips Curve (NKPC) has
co-moved with product-market \emph{structure}. Building on the \emph{Homothetic
Single Aggregator} (HSA) demand system of Matsuyama and Ushchev, we let the slope
depend on a low-frequency measure of the number of listed firms,
$\kt=\kappa_0+\delta\,\bar N_t$, and estimate the resulting state-space model by
Gibbs/FFBS as a theory-motivated, \emph{semi-structural} exercise. The coefficient
$\delta$ is robustly positive for the unemployment-gap specification (posterior mean
$+0.023$, $95\%$ credible interval $[+0.010,+0.039]$, Savage--Dickey Bayes factor
against $\delta=0$ of $\approx 79$, ``very strong''). Because the latent firm-count
trend $\bar N_t$ \emph{falls} over 1982--2012 (from about $+3.13$ to $-2.56$ in the
transformed units), a positive $\delta$ means the implied slope $\kt$ declines from
about $0.12$ to near zero \emph{as the number of firms falls---i.e.\ as market
concentration rises}. This sign is consistent with the HSA mechanism (fewer firms
$\Rightarrow$ higher markups and lower pass-through $\Rightarrow$ a flatter curve).
Unifying the left-hand side on \emph{core} CPI turns the HP output-gap specification
strongly positive as well ($\delta=+0.029$, Bayes factor $\approx 60$), while the
Beveridge--Nelson output gap stays inconclusive. The cross-measure sign disagreement
is concentrated in the 2008--09 episode and is consistent with an important role for
energy-price movements. Identification rests mainly on the 1980s, so we read the
result as \emph{theory-consistent time-series evidence, not a causal effect}.
\end{abstract}

\section*{Executive summary}
\addcontentsline{toc}{section}{Executive summary}

\paragraph{Question.} Has the slope of the NKPC co-moved with the number of
competing firms $N$? We estimate a semi-structural HSA-NKPC and read the sign and
magnitude of the coefficient $\delta$ in $\kt=\kappa_0+\delta\,\bar N_t$.

\paragraph{Direction, stated carefully.} $\delta>0$ means $\kt$ moves \emph{with}
$N$: the slope is steeper when firms are many and flatter when firms are few. Over
the sample the firm count $\bar N_t$ falls, so $\kt$ falls. The estimated
association is therefore ``\emph{fewer firms / rising concentration} $\to$
flattening,'' not ``rising competition $\to$ flattening.'' This is the direction the
HSA theory predicts.

\paragraph{Three findings.}
\begin{enumerate}
\item \textbf{The unemployment-gap slope flattens as the firm count falls.} With
core CPI, $\delta=+0.023$ $[+0.010,+0.039]$ and the Bayes factor against $\delta=0$
is $\approx 79$ (very strong). The implied $\kt$ declines from $\approx 0.12$ to
near zero as $\bar N_t$ falls.
\item \textbf{Unifying on core CPI aligns the output-gap evidence.} Under headline
CPI the HP output-gap slope shows no competition dependence; under core CPI it
becomes strongly positive ($\delta=+0.029$, Bayes factor $\approx 60$), matching the
unemployment result. The BN output gap stays inconclusive either way.
\item \textbf{The headline--core gap is concentrated in 2008--09.} A
Frisch--Waugh--Lovell / Okun decomposition places the sign disagreement almost
entirely in the six quarters 2008Q3--2009Q4, consistent with an important role for
energy prices; switching to core CPI removes most of it.
\end{enumerate}

\paragraph{Headline table.} Baseline priors, 1982Q1--2012Q4 ($T=124$), quarterly,
unrestricted, model \code{hsa\_steady}. Entries are posterior mean of $\delta$ with
$95\%$ credible interval; $\mathrm{BF}_{10}$ is the Savage--Dickey Bayes factor
\emph{for} $\delta\neq 0$ (concentration-dependent slope) against $\delta=0$.

\begin{table}[H]
\centering
\caption{Concentration dependence of the slope, headline vs.\ core CPI.
Recall $\delta>0$ with a falling $\bar N_t$ implies a flattening $\kt$.}
\begin{tabular}{lccc}
\toprule
Activity measure & Headline CPI: $\delta$ [95\% CI] ($\mathrm{BF}_{10}$)
& Core CPI: $\delta$ [95\% CI] ($\mathrm{BF}_{10}$)
& Core $\kt$ (start$\to$end) \\
\midrule
Unemployment gap & $+0.030\;[+0.009,+0.053]\;(20)$ & $+0.023\;[+0.010,+0.039]\;(79)$ & $0.120\to-0.005$ \\
Output gap (HP)  & $-0.002\;[-0.031,+0.026]\;(0.8)$ & $+0.029\;[+0.011,+0.049]\;(60)$ & $0.161\to-0.004$ \\
Output gap (BN)  & $-0.013\;[-0.041,+0.016]\;(1.2)$ & $+0.011\;[-0.009,+0.030]\;(0.9)$ & $0.062\to\phantom{-}0.005$ \\
\bottomrule
\end{tabular}
\end{table}

\noindent Bayes-factor reading (Kass--Raftery): $1$--$3$ anecdotal, $3$--$10$
moderate, $10$--$30$ strong, $30$--$100$ very strong, $>100$ decisive. The Bayes
factor tests the two-sided null $\delta=0$; the theory's one-sided prediction
$\delta>0$ is, if anything, favored more strongly.

\paragraph{Economic magnitude.} In the core-CPI unemployment specification the slope
falls from $0.120\;[0.065,0.176]$ at the start to
$-0.005\;[-0.045,0.035]$ at the end. The difference in the inflation response to a
one-percentage-point unemployment gap is therefore $0.125\;[0.056,0.195]$ percentage
point, and $0.502\;[0.224,0.781]$ for a four-point gap. These are model-implied
magnitudes, not causal contributions of the firm-count decline.

\paragraph{Scope condition.} The result is specific to samples that include the
1980s. On post-1988 samples the unemployment Phillips curve is already flat on
average, leaving little slope variation for any concentration trend to explain, and
$\delta\approx 0$. We report this openly (\autoref{sec:scope}).

\begin{figure}[H]
\centering
\includegraphics[width=0.82\linewidth]{delta_prior_posterior_by_spec.png}
\caption{Posterior distributions of $\delta$ across activity measures. The mass
shifts to the positive region for the unemployment and core output-gap
specifications.}
\end{figure}

\clearpage
\tableofcontents
\clearpage

\section{Introduction}
A central puzzle in inflation dynamics is the apparent decline over time in the
slope of the Phillips curve---the sensitivity of inflation to marginal cost or
slack---the so-called ``flattening.'' This paper asks whether that time variation
has co-moved with \emph{product-market structure}, specifically the low-frequency
trend in the number of firms, and evaluates the link quantitatively with a Bayesian
state-space model.

The theoretical basis is the \emph{Homothetic Single Aggregator} (HSA) demand
system of Matsuyama and Ushchev, as embedded in a New Keynesian model by Fujiwara
and Matsuyama (2026). In an HSA economy the strategic interaction among firms shapes
the curvature of the demand aggregator and hence the elasticity of price setting; as
a result the Phillips-curve slope $\kt$ depends on the number of competing firms.
Their central implication is that \emph{a fall in the number of firms}---rising
market concentration, a higher markup, and a lower pass-through---\emph{flattens} the
curve. Writing the low-frequency firm count as $\bar N_t$, this is a linear slope
equation
\begin{equation}
\kt=\kappa_0+\delta\,\bar N_t,
\end{equation}
with the theory predicting $\delta>0$: the slope co-moves with the number of firms
and therefore falls as concentration rises.

Our contribution is threefold, and we state it modestly. First, we provide, to our
knowledge, the first empirical operationalization of the HSA
concentration--Phillips-curve mechanism: we cast it as a state-space model---an
inflation observation equation together with a three-part decomposition of the firm
count---and estimate it as a \emph{theory-motivated, semi-structural} exercise
(\autoref{sec:distance} is explicit about the distance from the underlying theory).
Second, rather than only asking \emph{whether} the average slope flattened, we
estimate a \emph{driver} of the slope's time variation, separating a latent trend
and cycle in the firm count within the state space. Third, we anatomize where the
result comes from---which activity measures, which inflation measure, and which
sample quarters---and report the weak points openly.

\paragraph{Preview.} A positive $\delta$ (the slope flattens as the firm count
falls, i.e.\ as concentration rises) is robust for the unemployment gap; unifying on
core CPI aligns the output-gap evidence. The signal disappears on post-1988 samples,
which exclude the identifying 1980s variation, so we interpret the finding as
theory-consistent time-series evidence rather than a causal effect.

\section{Why we estimate this equation: the HSA Phillips curve}
\label{sec:why}
Our estimating equation is a theory-motivated counterpart of the New Keynesian
Phillips curve derived by Fujiwara and Matsuyama (2026, \emph{Journal of Monetary
Economics}) under a \emph{Homothetic Single Aggregator} (HSA) demand system. This
section states their result, maps it to what we take to the data, and is explicit
about how far the empirical specification sits from the theory.

\subsection{The HSA Phillips curve}
Final-goods producers assemble a continuum of varieties through an HSA demand
system with market-share function $s(z)$, where $z=p/A$ is a firm's price normalized
by the single price aggregator $A$ (the reference price for firms, an inverse
measure of competitive pressure). Intermediate-goods firms set prices under
Rotemberg adjustment costs $\chi$ and enter endogenously subject to a one-period
time-to-build lag. Solving the firm's problem, PPI inflation obeys (their Eq.~21)
\begin{equation}
\hat\pi_t=\beta(1-\delta)\,\mathbb E_t\hat\pi_{t+1}
+\underbrace{\frac{\zeta(z)-1}{\chi}}_{\text{slope on marginal cost}}
\big(\hat W_t-\hat Z_{P,t}-\hat p_t\big)
-\underbrace{\frac{1}{\chi}\,\frac{1-\rho(z)}{\rho(z)}}_{\text{entry / cost-push}}\hat N_t,
\label{eq:fm}
\end{equation}
where $\hat N_t$ is the (log) number of firms, $\hat W_t-\hat Z_{P,t}-\hat p_t$ is
the inverse markup (the model's real-marginal-cost measure),
$\zeta(z)=1-s'(z)z/s(z)>1$ is the price-elasticity function (so the flexible-price
markup is $\mu=\zeta/(\zeta-1)$), and $\rho(z)$ is the pass-through rate.

\subsection{Two ways competition enters}
\paragraph{(1) The slope is a function of the firm count (structural flattening).}
The coefficient on marginal cost, $(\zeta(z)-1)/\chi$, depends on the price
elasticity $\zeta(z)$, which through $z=p/A$ and the symmetric-equilibrium condition
$s(z)=1/N$ is a function of the number of firms $N$. Rising market
concentration---\emph{fewer firms}, a higher markup, a lower pass-through
$\rho(z)$---lowers this slope: the Phillips curve flattens \emph{structurally}.
Linearizing the slope in the low-frequency level of the firm count $\bar N_t$ gives
\begin{equation}
\kt=\kappa_0+\delta\,\bar N_t,
\end{equation}
so a \emph{positive} $\delta$ is the theory's flattening prediction: the slope
co-moves with the number of firms, i.e.\ it falls as concentration rises and $N$
falls.

\paragraph{(2) Endogenous entry is a cost-push term (estimation bias).} The second
term, $-\tfrac{1}{\chi}\tfrac{1-\rho(z)}{\rho(z)}\hat N_t$, is present only away from
CES. It says the cyclical component of firm entry acts as a cost-push shock through
strategic complementarity; omitting it biases empirical Phillips-curve estimates (an
omitted-variable problem). We keep the cyclical term explicitly, with its own
coefficient $\theta_t=\theta_0+\gamma\,\bar N_t$, but note below that our workhorse
sets it to zero.

\paragraph{CES is the irrelevance case.} Under CES,
$s(z)=\gamma_{\mathrm{CES}}\,z^{1-\theta}$, so $\zeta(z)=\theta$ is constant and the
pass-through is $\rho(z)=1$. The entry term vanishes and the slope is the constant
$(\theta-1)/\chi$: market concentration neither flattens the curve nor shifts
inflation. This is exactly our \code{ces} benchmark, and it is why departing from CES
(our HSA models) is what allows $\delta$ and the $\hat N_t$ term to be nonzero in the
first place.

\subsection{From theory to the estimating equation}
Writing the activity/inverse-markup gap as $x_t$, adding a backward-looking
indexation term ($\alpha$) and a reduced-form shock $e_t$, and decomposing the
observed firm count into a trend and a cycle,
\begin{equation}
N^{\mathrm{obs}}_t=\bar N_t+\hat N_t+\nu_t,
\end{equation}
with $\bar N_t$ a random-walk-with-drift trend, $\hat N_t$ a stationary AR(2) gap
(truncated to the stationary region), and $\nu_t\sim\mathcal N(0,\sigma_N^2)$
measurement error, the HSA Phillips curve becomes the equation we estimate:
\begin{equation}
\boxed{\;\pi_t=\alpha\,\pi_{t-1}+(1-\alpha)\,\mathbb E_t\pi_{t+1}
+\kt\,x_t-\theta_t\,\hat N_t+e_t,\qquad
\kt=\kappa_0+\delta\,\bar N_t,\;\;\theta_t=\theta_0+\gamma\,\bar N_t.\;}
\label{eq:est}
\end{equation}
The trend $\bar N_t$ carries the structural-flattening channel (through $\delta$);
the cycle $\hat N_t$ carries the entry cost-push channel (through $\theta_t$).

\subsection{Distance from the theory}
\label{sec:distance}
We call the exercise \emph{semi-structural} rather than structural because the
specification departs from \eqref{eq:fm} in several ways, each of which we flag so
the coefficients are not over-interpreted:
\begin{itemize}
\item the theory is about \emph{PPI} inflation, whereas we estimate on CPI / core CPI;
\item the theoretical regressor is real marginal cost / the inverse markup, whereas
our main results use the unemployment gap and the HP output gap;
\item we do not impose the theory's cross-coefficient restrictions, and approximate
$\kt$ as a \emph{linear} function of $\bar N_t$;
\item we add backward indexation ($\alpha$), which is not in \eqref{eq:fm};
\item expected inflation is an external survey measure (see \autoref{sec:data});
\item our workhorse \code{hsa\_steady} sets the cyclical entry term $\theta_t$ to
zero, i.e.\ it uses only the trend channel.
\end{itemize}
Accordingly, testing $\delta>0$ is best read as a \emph{reduced-form test consistent
with the sign restriction} of the Fujiwara--Matsuyama mechanism, not as a structural
estimate of its deep parameters. The mechanism implies (i) $\delta>0$ and (ii) a
monotone decline of $\kt$ as $\bar N_t$ falls; $\delta=0$ (with $\theta=0$) collapses
\eqref{eq:est} to the CES benchmark---exactly the constant- vs.\ time-varying-slope
comparison of \autoref{sec:bench}. Sections~\ref{sec:main}--\ref{sec:bench} examine
both.

\section{Data}
\label{sec:data}
\paragraph{Sample.} Quarterly, 1982Q1--2012Q4 ($T=124$), with robustness
sub-samples. Complete-case availability, rather than the wider date range of the
merged source file, determines each estimation window: pre-2008 contains 103
quarters (1982Q1--2007Q3) and start-1988 contains 100 quarters
(1988Q1--2012Q4). The nominal post-2008 window contains only 16 complete quarters
(2009Q1--2012Q4), below our 40-observation minimum, and is therefore not treated as
an informative robustness exercise.

\paragraph{Inflation.} Headline CPI (\code{pi\_cpi}) and \textbf{core} CPI
(\code{pi\_cpi\_core}); core PCE as an auxiliary check. Expected inflation
$\mathbb{E}\pi$ is an external survey measure held on the \emph{headline} basis
across inflation specifications. This introduces a headline--core expectations
mismatch that we discuss as a limitation in \autoref{sec:ident} and
\autoref{sec:status}: the headline--core wedge is time-varying and correlates with
energy prices and the cycle, so treating it as absorbed by the constant and $\alpha$
is only an approximation, and the sensitivity of the results to headline vs.\ core is
itself a symptom of it.

\paragraph{Activity gaps (five).} Unemployment gap; BN output gap; HP output gap
($100\times$ log real output, HP-filtered, same $100$-log-point unit as BN);
HP labor-share gap; inverse markup.

\paragraph{Competition proxy.} The competition series is the Gustavo firm-count
measure (\code{N\_Gustavo}, annual, PCHIP-interpolated to quarterly): the total
number of U.S.\ listed firms. The theoretical $N$ is the number of firms competing
\emph{strategically within a market}, which is not the same object. We therefore use
\code{N\_Gustavo} as an \emph{aggregate proxy for the mass of active/listed firms,
intended to approximate the theoretical number of competitors}. The two can diverge:
unlisted firms may grow while listed firms shrink, industry composition changes,
merger and delisting rules change, and foreign or import competition can intensify;
moreover an economy-wide firm count is not the same as within-industry
concentration. We make this measurement scope explicit and do not equate the total
listed-firm count with generic ``competition'' or within-industry concentration.
Over the sample the transformed latent trend $\bar N_t$ falls from about $+3.13$ to $-2.56$
(\autoref{fig:data}c), which is why $\delta>0$ delivers a \emph{falling} $\kt$.

\paragraph{Unit convention.} The default transform is
$(100\log N-\overline{100\log N})/10$, so one unit of $\bar N$ or $\hat N$ is a
ten-log-point move around the sample mean, and reported $\delta,\theta,\theta_0,\gamma$
are all on this ten-log-point scale.

\begin{figure}[H]
\centering
\includegraphics[width=0.92\linewidth]{report/data_series.png}
\caption{Data. (a) Headline vs.\ core CPI inflation; the shaded band marks the
2008--09 financial crisis, during which the collapse in oil prices drove headline
CPI sharply negative (to about $-1.6\%$) while core CPI stayed near $1.5\%$. (b)
Real-activity gaps. (c) The listed-firm-count proxy under the estimation transform;
its low-frequency trend \emph{falls} over the sample and is the driver of the slope
equation.}
\label{fig:data}
\end{figure}

\section{Econometric model and estimation}
\label{sec:model}

\subsection{Model variants}
Five models differ only in the slope restrictions (\autoref{tab:models}).

\begin{table}[H]
\centering
\caption{Model variants.}
\label{tab:models}
\begin{tabular}{llll}
\toprule
Model & $\kt$ & $\theta_t$ & Role \\
\midrule
\code{ces}             & constant & $0$ & textbook NKPC benchmark \\
\code{hsa\_steady}     & $\kappa_0+\delta\bar N_t$ & $0$ & \textbf{workhorse}: trend channel only \\
\code{hsa\_dynamic}    & constant & constant & $\hat N$ channel only \\
\code{hsa\_const\_theta} & $\kappa_0+\delta\bar N_t$ & constant & \code{hsa\_full} with $\gamma=0$ \\
\code{hsa\_full}       & $\kappa_0+\delta\bar N_t$ & $\theta_0+\gamma\bar N_t$ & full model \\
\bottomrule
\end{tabular}
\end{table}

\noindent The workhorse \code{hsa\_steady} uses only the trend channel ($\theta_t=0$):
it is the cleanest place to read $\delta$, at the cost of dropping the theory's
cyclical entry term. \code{hsa\_const\_theta} exists because the $\gamma$ regressor
$\hat N\!\cdot\!\bar N$ is nearly collinear (about $98\%$) with the $\theta_0$
regressor $\hat N$, so $\gamma$ is not identified.

\subsection{Estimation}
The three firm-count components are drawn jointly by Kalman/FFBS; coefficient blocks
are updated from their conjugate Gaussian conditionals (with rejection sampling when
sign constraints are imposed). \code{hsa\_full} carries an explicit N-measurement
error $\sigma_N^2$, and its two-block state sampler ($\hat N\mid\bar N$, then
$\bar N\mid\hat N$) is an exact Gibbs step. Kappa priors are in physical units
(converted internally by $\mathrm{KAPPA\_SCALE}=100$); N-shock and
measurement-variance priors are in squared ten-log-point units ($\approx 0.01$
decade). The workhorse \code{hsa\_steady} converges well ($\hat R\le 1.01$); some
\code{hsa\_full} cells (weak/tight/annual-Q4) show $\hat R\in[1.1,1.3]$ and are not
over-interpreted.

\begin{figure}[H]
\centering
\includegraphics[width=0.72\linewidth]{unemployment_gap_core/competition_decomposition_path.png}
\caption{Posterior firm-count decomposition $N=\bar N+\hat N$ (unemployment gap,
core CPI). Trend and gap are well identified (RMSE $\le 0.01$) in every
specification; clean states, however, do not by themselves identify $\kappa$ or
$\delta$.}
\end{figure}

\section{Main results}
\label{sec:main}

\subsection{The unemployment-gap slope flattens as the firm count falls}
For the baseline \code{hsa\_steady} on the main sample, the core-CPI unemployment
specification gives
\[
\delta=+0.0227\;[+0.0099,+0.0393],\qquad \kappa_0=+0.051\;[+0.015,+0.088],
\]
with a Savage--Dickey Bayes factor against $\delta=0$ of $\approx 79$ (very strong).
Because $\bar N_t$ falls over the sample, the implied slope $\kt$ declines
monotonically from $0.120$ to about zero: as the firm count falls (concentration
rises), the Phillips curve flattens, the sign the HSA mechanism predicts. Under
headline CPI the same coefficient is $+0.030\;[+0.009,+0.053]$ (Bayes factor
$\approx 20$, strong), so the unemployment result is robust to the inflation measure.

\paragraph{On ``near zero.''} The endpoint posterior mean of $\kt$ is marginally
\emph{negative} ($-0.005$ core; $-0.010$ headline), and the endpoint credible band
includes zero (\autoref{fig:kt_unemp}). ``Flattens to near zero'' is therefore the
accurate statement, not a strictly positive residual slope; a $\kt\ge0$-constrained
version (\autoref{sec:robust}) delivers the same flattening without the small
negative overshoot.

\begin{figure}[H]
\centering
\includegraphics[width=0.82\linewidth]{report/kappa_t_unemployment.png}
\caption{Time-varying slope $\kt$ for the unemployment gap (posterior mean and
$95\%$ band), headline and core CPI. Both decline to about zero by the end of the
sample; the endpoint band includes zero.}
\label{fig:kt_unemp}
\end{figure}

\subsection{Economic magnitude: how many inflation points does flattening imply?}
\label{sec:magnitude}
We translate statistical identification into economic units. In the core-CPI
unemployment specification the start-to-end slope difference is
$0.125\;[0.056,0.195]$. Thus, for an unemployment-gap movement of $g$ percentage
points, the difference between the two dates' inflation responses is
$g(\kappa_{\mathrm{start}}-\kappa_{\mathrm{end}})$ percentage points.

\begin{table}[H]
\centering
\caption{Economic magnitude of the time-varying slope (core CPI, unemployment gap).
Each row is the difference between the 1982 and 2012 inflation responses to the same
unemployment-gap movement.}
\label{tab:magnitude}
\begin{tabular}{lcc}
\toprule
Unemployment-gap movement & Posterior mean response difference & $95\%$ credible interval \\
\midrule
1 percentage point & 0.125 & [0.056, 0.195] \\
2 percentage points & 0.251 & [0.112, 0.391] \\
4 percentage points & 0.502 & [0.224, 0.781] \\
One sample standard deviation (1.65 points) & 0.207 & [0.092, 0.322] \\
\bottomrule
\end{tabular}
\end{table}

\noindent A four-point deterioration ($x_t=-4$), for example, lowers core inflation by
about 0.48 percentage point under the start-of-sample slope but has almost no effect
under the end-of-sample slope. Along the realized unemployment-gap path, the
difference between holding the slope fixed at its initial value and using the
time-varying slope has an RMS of $0.176\;[0.078,0.275]$ point and a maximum absolute
value of $0.616\;[0.272,0.964]$ point. This is a within-model counterfactual, not a
causal decomposition attributing the difference to the firm-count decline.

\subsection{Unifying on core CPI aligns the output-gap evidence}
\label{sec:core}
For the output-gap specifications the concentration dependence is absent or negative
under headline CPI but turns positive under core CPI (\autoref{tab:coreflip}). The
HP output gap in particular becomes strongly positive ($\delta=+0.029$, Bayes factor
$\approx 60$), matching the unemployment result. The BN output gap remains
inconclusive under both inflation measures (Bayes factors $\approx 1$; the credible
interval spans zero), so the corroborating evidence comes from HP, not BN. Reading
the credible intervals honestly, ``directionally positive but weak'' is the right
description of BN core, and the earlier impression of BN as supportive evidence does
not survive a Bayes-factor assessment.

\begin{table}[H]
\centering
\caption{Concentration dependence $\delta$ of the slope, headline vs.\ core CPI
(\code{hsa\_steady}, baseline). $\mathrm{BF}_{10}$ is the Savage--Dickey Bayes
factor for $\delta\neq 0$.}
\label{tab:coreflip}
\begin{tabular}{lccc}
\toprule
Output gap & Headline $\delta$ [95\% CI] & Core $\delta$ [95\% CI] & $\mathrm{BF}_{10}$ (head / core) \\
\midrule
BN & $-0.0128\;[-0.0411,+0.0161]$ & $+0.0106\;[-0.0092,+0.0301]$ & $1.2\;/\;0.9$ \\
HP & $-0.0024\;[-0.0312,+0.0260]$ & $+0.0289\;[+0.0108,+0.0494]$ & $0.8\;/\;60$ \\
(Unemployment, ref.) & $+0.0302\;[+0.0092,+0.0527]$ & $+0.0227\;[+0.0099,+0.0393]$ & $20\;/\;79$ \\
\bottomrule
\end{tabular}
\end{table}

\noindent After unifying on core CPI, $\kt$ declines over the sample for the
unemployment, HP, and (mildly) BN specifications, whereas under headline CPI the BN
slope perversely \emph{steepens} ($0.025\to0.096$)---an artifact we discuss in
\autoref{sec:ident}. The same pattern holds under \code{hsa\_full}, so the finding
is robust across model variants.

\begin{figure}[H]
\centering
\includegraphics[width=0.86\linewidth]{report/kappa_t_others.png}
\caption{Time-varying slope $\kt$ for the non-unemployment activity measures
(posterior means). Core-CPI HP (dashed purple) flattens sharply; headline BN
(solid green) steepens---the episode we examine in \autoref{sec:ident}.}
\label{fig:kt_others}
\end{figure}

\subsection{Slope level $\kappa_0$: is the curve alive?}
The slope at average competition $\kappa_0$ is positive for the unemployment and
output-gap specifications: HP headline $+0.142\;[+0.048,+0.233]$, HP core
$+0.070\;[+0.017,+0.124]$, unemployment core $+0.051\;[+0.015,+0.088]$; the credible
interval excludes zero in each case. Switching to core lowers the level (the
headline value was inflated by energy prices) but preserves the positive sign. For
the labor-share gap and the inverse markup---the latter being the theory's most
direct regressor---the intervals are wide and include zero: these specifications
carry essentially no inflation signal, and their $\kappa_0$ and $\delta$ are
prior-dominated. That the result is weakest for the inverse markup is itself a
caution about the mapping to the theory.

\subsection{Consistency of the intermediate channel: firm counts and aggregate markups}
\label{sec:markup_mechanism}
The HSA intermediate channel runs from fewer firms through higher markups and lower
pass-through to a flatter slope. Over the same sample the listed-firm count falls
$21.2\%$, from $7.57$ in 1982 to $5.97$ in 2012; the aggregate markup rises
$12.6\%$, from $1.580$ to $1.779$; and the inverse markup falls $11.2\%$. The
level correlation between the firm count and markup is $-0.63$, in the theory's
direction, but the first-difference correlation is $+0.05$ and the linearly
detrended correlation only $-0.20$. This is consistency in common long-run trends,
not mechanism identification from independent short-run movements.

\begin{table}[H]
\centering
\caption{Aggregate diagnostic of the HSA intermediate channel. Endpoints are
four-quarter averages within each year.}
\label{tab:markup_mechanism}
\begin{tabular}{lrrr}
\toprule
Series & 1982 & 2012 & Percent change \\
\midrule
Listed-firm count & 7.574 & 5.971 & $-21.2\%$ \\
Aggregate markup & 1.580 & 1.779 & $+12.6\%$ \\
Inverse markup & 0.633 & 0.562 & $-11.2\%$ \\
\bottomrule
\end{tabular}
\end{table}

\noindent Moreover, with $\mu=\zeta/(\zeta-1)$,
$\kappa=(\zeta-1)/\chi$, and fixed $\chi$, we have
$\kappa\propto1/(\mu-1)$. Starting from the estimated 1982 slope of $0.120$ and
moving only the observed markup predicts a 2012 slope of $0.089$. This corresponds
to about $24.5\%$ of the estimated $0.125$ flattening. The direction is HSA-consistent,
but the observed markup alone does not account for most of the decline to near zero.
Because the aggregate markup is not identical to the model's structural markup and
fixed $\chi$ is a strong assumption, this is a calibration-style magnitude check.

\section{When the slope is constant: benchmark and the case for time variation}
\label{sec:bench}
To evaluate the time-varying result of \autoref{sec:main} we need a constant-slope
benchmark. This section presents the constant-slope models (\code{ces} and
\code{hsa\_dynamic}) and then asks whether letting the slope vary with the firm
count improves in-sample fit.

\subsection{Constant-slope models (\code{ces} / \code{hsa\_dynamic})}
\paragraph{CES (constant $\kappa$).} The textbook NKPC slope is clearly positive for
some activity measures (\autoref{tab:ceskappa}): the HP output gap and the core-CPI
unemployment gap have credible intervals that exclude zero. The labor-share gap and
inverse markup cannot be distinguished from zero even as a constant, confirming the
absence of an inflation signal there.

\begin{table}[H]
\centering
\caption{Constant slope $\kappa$ (\code{ces}, baseline): posterior mean [95\% CI].}
\label{tab:ceskappa}
\begin{tabular}{lcc}
\toprule
Activity measure & Headline CPI & Core CPI \\
\midrule
Unemployment gap    & $+0.056\;[-0.008,+0.120]$ & $+0.039\;[+0.007,+0.072]$ \\
Output gap (HP)     & $+0.138\;[+0.047,+0.227]$ & $+0.091\;[+0.043,+0.139]$ \\
Output gap (BN)     & $+0.059\;[-0.025,+0.144]$ & $+0.032\;[-0.017,+0.081]$ \\
Labor-share gap (HP)& $-0.008\;[-0.169,+0.150]$ & --- \\
Inverse markup      & $+0.007\;[-0.217,+0.228]$ & --- \\
\bottomrule
\end{tabular}
\end{table}

\begin{figure}[H]
\centering
\includegraphics[width=0.82\linewidth]{report/pp_kappa_ces_vs_steady.png}
\caption{Prior vs.\ posterior of the Phillips-curve slope for the unemployment gap:
CES's constant slope $\kappa$ and \code{hsa\_steady}'s slope at average competition
$\kappa_0$, against their common prior $\mathcal N(0.1,0.2)$. Both are sharply
updated by the data and sit just above zero; the two models agree on the slope
level, so what distinguishes them is whether that slope varies with the firm count
(\autoref{sec:mc}).}
\label{fig:kappa_pp}
\end{figure}

\paragraph{\code{hsa\_dynamic} (constant $\kappa$ and $\theta$).} Adding the
firm-count-gap channel $\hat N$ leaves the constant slope essentially unchanged from
CES. The added constant $\theta$ is weakly negative but effectively unidentified at
quarterly frequency: its $95\%$ interval includes zero (unemployment
$-0.019\;[-0.146,+0.153]$; HP $-0.017\;[-0.122,+0.123]$; BN $-0.069\;[-0.165,+0.094]$).
The quarterly variation in $\hat N$ is an artifact of annual-to-quarterly
interpolation, and at annual (Q4) frequency $\theta$ reverts to its prior. The
$\hat N$ channel therefore carries little information in these data.

\begin{figure}[H]
\centering

\includegraphics[width=0.66\linewidth]{prior_posterior_theta.png}
\caption{Prior vs.\ posterior of $\theta$. The posterior does not separate from the
prior: $\theta$ is not identified at quarterly frequency.}
\end{figure}

\subsection{Does the time-varying slope improve fit over the constant slope?}
\label{sec:mc}
We compare CES (constant $\kappa$) against \code{hsa\_steady} (slope varying with the
firm count) by two criteria: an in-sample fit score and the Savage--Dickey Bayes
factor for $\delta\neq 0$ (\autoref{tab:mc}).

\paragraph{In-sample plug-in fit score.} We evaluate the inflation equation at
posterior-mean parameters and states (bars denote posterior means),
\begin{equation}
\hat\pi_t=\bar\alpha\,\pi_{t-1}+(1-\bar\alpha)\,\mathbb{E}_t\pi_{t+1}+\bar\kt\,x_t,
\qquad \bar\kt=\bar\kappa_0+\bar\delta\,\bar N_t
\;\;(\text{CES: }\bar\kt=\bar\kappa),
\end{equation}
and score the fit by the Gaussian log-likelihood of the residuals
$r_t=\pi_t-\hat\pi_t$ with $\bar\sigma_e^2$ the posterior-mean shock variance:
\begin{equation}
S=-\frac{T}{2}\log\!\big(2\pi\bar\sigma_e^2\big)-\frac{1}{2\bar\sigma_e^2}\sum_{t=1}^{T}r_t^2 .
\end{equation}
This is an \emph{in-sample plug-in} score: it reuses the estimation sample and the
posterior means, and is neither an out-of-sample forecast nor a posterior predictive
density. It is used only as a rough comparative fit diagnostic; a proper predictive
comparison (LOO, WAIC, rolling one-step forecasts, or an integrated log predictive
density) is left for future work.

\paragraph{Bayes factor (Savage--Dickey), stated precisely.} We test $\delta=0$
\emph{within the HSA state model}: $M_1$ is \code{hsa\_steady} with $\delta$ free and
$M_0$ is the same model with $\delta=0$. The Savage--Dickey identity gives
\begin{equation}
\mathrm{BF}_{10}=\frac{p(\delta=0)}{\pi(\delta=0\mid y)},\qquad
p(\delta=0)=\frac{1}{\sigma_\delta\sqrt{2\pi}}=\frac{1}{0.02\sqrt{2\pi}}\approx 19.95
\;\;(\delta\sim\mathcal N(0,\sigma_\delta^2)),
\end{equation}
with the posterior ordinate estimated by a Gaussian kernel density on the draws
$\{\delta^{(s)}\}$, $\hat\pi(\delta=0\mid y)=\tfrac{1}{Sh}\sum_s\phi\!\big((0-\delta^{(s)})/h\big)$.
Three cautions. (i) This is \emph{not} a CES-vs-HSA marginal-likelihood ratio: CES
and HSA also differ in the latent-$N$ block, which the SDDR holds fixed, so we keep
the comparison inside the HSA state model. (ii) The theory is one-sided ($\delta>0$);
the two-sided $\mathrm{BF}_{10}$ is conservative. (iii) The zero-point posterior
density is a KDE estimate; the ordering across parameters is insensitive to the
bandwidth, but the exact magnitudes are not, so we report them to one or two
significant figures.

\begin{table}[H]
\centering
\caption{Constant vs.\ time-varying slope (baseline, quarterly, unrestricted).
The fit score is the in-sample plug-in $S$ above (higher is better);
$\mathrm{BF}_{10}$ is the SDDR for $\delta\neq0$ within the HSA state model.}
\label{tab:mc}
\begin{tabular}{lccc}
\toprule
Activity measure & CES fit score & \code{hsa\_steady} fit score & $\mathrm{BF}_{10}(\delta\neq0)$ \\
\midrule
Unemployment gap     & $-119.9$ & $\mathbf{-115.3}$ & $\mathbf{20.2}$ \\
Output gap (BN)      & $-120.8$ & $-119.6$ & $1.2$ \\
Output gap (HP)      & $-114.1$ & $-114.0$ & $0.8$ \\
Labor-share gap (HP) & $-121.9$ & $-121.4$ & $0.9$ \\
Inverse markup       & $-121.9$ & $-121.9$ & $1.0$ \\
\bottomrule
\end{tabular}
\end{table}

\paragraph{Verdict.} For the unemployment gap, allowing the slope to vary with the
firm count improves the in-sample fit score by about $4.6$ and the Bayes factor
favors $\delta\neq0$ by roughly $20$:$1$. For the headline output gaps the gain is
marginal, consistent with Sections~\ref{sec:core}--\ref{sec:ident}: the output-gap
signal lives in the core-CPI specifications, which are not yet in the
model-comparison pipeline (see \autoref{sec:status}). For the labor-share gap and
inverse markup there is no inflation signal, so constant and time-varying slopes fit
equally well. We stress that this is an in-sample fit comparison, not evidence of
out-of-sample predictive gains.

\paragraph{On marginal likelihoods.} This verdict rests on the fit score and the
Savage--Dickey ratios. The Chib marginal likelihood is contaminated for annual-Q4
blocks by a missing-N bug (\autoref{sec:status}) and is not used.

\subsection{Savage--Dickey ratios, parameter by parameter}
\autoref{tab:sddr} reports the Savage--Dickey Bayes factor for the point null
``$\theta=0$'' for every scalar coefficient of every model (unemployment gap,
baseline). Entries are $\mathrm{BF}_{10}=p(\theta=0)/\hat\pi(\theta=0\mid y)$
(the reciprocal of the Savage--Dickey density ratio); larger favors ``$\ne 0$'', and
$<1$ favors the null. A dot means the parameter is absent from that model.

\begin{table}[H]
\centering
\caption{Savage--Dickey Bayes factors $\mathrm{BF}_{10}$ for ``parameter $=0$'', by
model (unemployment gap, headline CPI, baseline). Top block: coefficients of
economic interest. Bottom block: auxiliary dynamics, for which the null $=0$ is
mechanical (these parameters are obviously nonzero).}
\label{tab:sddr}
\begin{tabular}{lccccc}
\toprule
Parameter & \code{ces} & \code{hsa\_dyn} & \code{hsa\_steady} & \code{hsa\_c\_theta} & \code{hsa\_full} \\
\midrule
$\kappa$ (const.\ slope)   & $0.61$ & $0.39$ & $\cdot$ & $\cdot$ & $\cdot$ \\
$\kappa_0$ (slope level)   & $\cdot$ & $\cdot$ & $1.07$ & $3.76$ & $1.95$ \\
$\delta$ (comp.\ dependence) & $\cdot$ & $\cdot$ & $\mathbf{20.2}$ & $\mathbf{13.0}$ & $\mathbf{8.9}$ \\
$\theta$ ($\hat N$ channel)  & $\cdot$ & $0.49$ & $\cdot$ & $0.50$ & $\cdot$ \\
$\theta_0$                 & $\cdot$ & $\cdot$ & $\cdot$ & $\cdot$ & $0.47$ \\
$\gamma$                   & $\cdot$ & $\cdot$ & $\cdot$ & $\cdot$ & $1.35$ \\
\midrule
$\alpha$ (indexation)      & $>10^4$ & $>10^4$ & $>10^4$ & $>10^4$ & $>10^4$ \\
$\phi_1$                   & $>10^4$ & $>10^4$ & $>10^4$ & $>10^4$ & $>10^4$ \\
$\rho_1$ (AR2)             & $\cdot$ & $>10^4$ & $>10^4$ & $>10^4$ & $>10^4$ \\
$\rho_2$ (AR2)             & $\cdot$ & $>10^4$ & $>10^4$ & $>10^4$ & $>10^4$ \\
$n$ (drift)                & $\cdot$ & $0.83$ & $32.5$ & $>10^4$ & $>10^4$ \\
\bottomrule
\end{tabular}
\end{table}

\noindent The reading matches the rest of the paper: $\delta$ is the one
economically interesting coefficient with clear evidence against zero in every HSA
model, whereas the constant slope ($\kappa,\kappa_0$) is barely distinguishable from
zero and the $\hat N$-channel coefficients ($\theta,\theta_0,\gamma$) are not
identified ($\mathrm{BF}_{10}\lesssim 1$). Variance parameters are omitted because
the null ``standard deviation $=0$'' is degenerate under the Savage--Dickey ratio.

\section{Where the identifying variation sits}
\label{sec:ident}
\subsection{Can the firm-count trend be distinguished from a generic time trend?}
The posterior mean $\bar N_t$ in the main specification is nearly linear in time:
its correlation with a linear time trend is $-0.99994$. The correlation between the
slope interaction $x_t\bar N_t$ and $x_t\times\mathrm{time}_t$ is $-0.99455$ and,
after FWL-removing the inertia term and $x_t$, becomes $-0.99999$. In a horse race
containing both interactions, the $t$ statistic is $1.16$ on the firm-count
interaction and $1.13$ on the time interaction. The design detects a slope decline
that co-moves with the firm count, but cannot distinguish it from a generic linear
change in the slope.

\begin{figure}[H]
\centering
\includegraphics[width=0.96\linewidth]{report/competition_trend_identification.png}
\caption{Firm-count trends and identifying information. (a) Observed listed-firm
count, source-data BN trend, and posterior latent trend in the same ten-log-point
units. The observed series and source BN trend rise until the mid-1990s, while the
posterior latent trend falls nearly linearly. (b) Each period's share of the squared
FWL-residualized $x_t\bar N_t$ interaction.}
\label{fig:trend_identification}
\end{figure}

\paragraph{Sensitivity to the trend decomposition.}
The transformed observed count moves from $-0.23$ at the start to a peak of $+2.16$
around 1996 and then to $-2.69$ at the end. The posterior $\bar N_t$, by contrast,
falls almost monotonically from $+3.13$ to $-2.56$. Its start-to-end decline is
$-5.68$ (about 56.8 log points), compared with $-2.47$ (24.7 log points) in the
observed series. OLS diagnostics using the same core-CPI unemployment equation retain
the sign under alternative trends but cut the initial slope roughly in half.

\begin{table}[H]
\centering
\small
\caption{Sensitivity to the firm-count trend decomposition. The posterior-$\bar N$
row is an OLS analog using the main model's posterior mean as a regressor; none of the
rows is a Bayesian re-estimation. Standard errors are in parentheses.}
\label{tab:trend_sensitivity}
\begin{tabular}{lrrrr}
\toprule
Firm-count series used in slope & Series change & $\delta$ & $\kappa_{\mathrm{start}}$ & $\kappa_{\mathrm{end}}$ \\
\midrule
Posterior latent trend $\bar N$ & $-5.683$ & $0.025\;(0.006)$ & 0.126 & $-0.017$ \\
Observed firm count & $-2.467$ & $0.022\;(0.010)$ & 0.060 & \phantom{-}0.005 \\
Source-data BN trend & $-2.665$ & $0.024\;(0.010)$ & 0.061 & $-0.003$ \\
\bottomrule
\end{tabular}
\end{table}

\paragraph{Which periods identify the interaction?}
Of the residual interaction variation, $60.1\%$ is in 1982--87 and $34.5\%$ in
2008--12; 1988--99 contributes $4.6\%$ and 2000--07 only $0.9\%$. It is more
accurate to say that \emph{the two ends of the sample---the early 1980s and the
post-crisis years---identify the result, while the middle two decades contribute
almost nothing} than to attribute identification to the 1980s alone.

\subsection{Why do headline and core differ?}
The cross-measure sign disagreement is data information concentrated in a few
quarters rather than a uniform feature of the sample. An OLS-analog
Frisch--Waugh--Lovell (FWL) and Okun decomposition ($T=124$, own-specification
posterior $\bar N$) shows the following.

\begin{enumerate}
\item \textbf{Episode contribution.} Of the BN total $\delta=-0.046$, the six
quarters 2008Q3--2009Q4 account for $-0.046$ (about $100\%$). Dropping them flips BN
to $+0.088$ (t $\approx 2.1$), the same sign as unemployment ($+0.058$). The
1983--86 Volcker period contributes only $-0.003$.
\item \textbf{Okun decomposition.} The BN gap shares only $R^2=0.11$ of its variance
with the unemployment gap. Its Okun-common component of $\delta$ is $+0.026$
(matching the unemployment posterior $+0.029$); the Okun-\emph{residual} component
$-0.072$ carries the entire disagreement.
\item \textbf{Mechanism, stated cautiously.} In 2008--09 output collapsed earlier and
deeper than the labor market (Okun residual $-2.4$) exactly when energy prices swung
headline CPI from $+2.3$ to $-2.2$. BN reads this as ``huge slack $\times$ huge
disinflation $=$ steep slope,'' and in the lowest-$\bar N$ era, giving $\delta<0$.
\item \textbf{Core CPI moves the measures together.} Replacing the left-hand side
with core, or dropping the six quarters, moves BN, HP, and unemployment to the
positive side (\autoref{sec:core}).
\end{enumerate}

We resist a strong ``it is entirely an oil artifact'' claim. Dropping the six
quarters, or switching to core, simultaneously removes energy-price movements
\emph{and} crisis-era measurement error, the sharp GDP--unemployment divergence, and
financial-constraint effects; and the expectations term stays headline-based
throughout. The defensible statement is: \emph{the sign disagreement is concentrated
in the 2008--09 episode and is consistent with an important role for energy-price
movements}. As corroboration, on the available pre-2008 sample
(1982Q1--2007Q3) the BN headline
$\delta$ moves from $-0.013$ to $+0.002$: the negative sign vanishes once the episode
is excluded.

\section{Robustness}
\label{sec:robust}
We separate two questions: does $\delta$ keep a positive \emph{sign}, and is it
\emph{identified away from zero}? Under weak and tight priors the unemployment
$\delta$ stays in $[+0.02,+0.04]$ with credible intervals that exclude zero. At
annual (Q4) frequency---N observed only in Q4, Q1--Q3 treated as missing and inferred
by the state-space model---the sign is unchanged; the quarterly ``significance'' of
$\theta/\gamma$, by contrast, is an interpolation artifact that reverts to the prior
at annual frequency. In the usable period comparisons the sign is stable, but
identification is not uniform: once the 1980s are excluded in the 100-quarter
start-1988 sample, the interval widens toward zero even though the mean stays
positive (\autoref{sec:scope}). We do not interpret the 16-quarter post-2008
estimate. Imposing the hard constraint $\kt\ge0$ (a different, truncated posterior,
reported separately) preserves the sign.

\paragraph{Time-trend diagnostic.} The linear-time-trend horse race in
\autoref{sec:ident} shows that $x_t\bar N_t$ and $x_t\times\mathrm{time}_t$ are
effectively indistinguishable. Quadratic trends, structural-break dummies, and placebo
interactions with unrelated trending variables remain untested. Thus $\delta$ should
be read as an association, not a causal effect (\autoref{sec:scope}).

\begin{figure}[H]
\centering
\includegraphics[width=0.72\linewidth]{prior_sensitivity_delta.png}
\caption{Prior sensitivity of $\delta$. For the unemployment specification the
posterior is not driven by the prior; the data pin it in the positive region.}
\end{figure}

\section{Scope condition and the limits of identification}
\label{sec:scope}
For honesty we state where the main result does \emph{not} hold and what it can and
cannot claim. The identifying variation is concentrated in the 1980s, so the sample
period matters.
\begin{itemize}
\item \textbf{Post-1988 samples.} Restricting the start date to 1988 weakens the
unemployment $\delta$ (its credible interval moves toward zero): after 1988 the
unemployment Phillips curve is already flat on average ($\kappa_0\approx0$), leaving
little slope variation for any firm-count trend to explain.
\item \textbf{The two sample endpoints are the engine.} Residual interaction variation
is concentrated in 1982--87 ($60.1\%$) and 2008--12 ($34.5\%$). The former combines
the sharp Volcker disinflation with a high \emph{latent} firm-count trend; the latter
combines large post-crisis slack with a low latent trend. The observed firm count
itself peaks in the mid-1990s, not in the 1980s.
\end{itemize}

\paragraph{Not a causal effect.} $\bar N_t$ is a near-monotone low-frequency trend,
so $\delta$ cannot be separated, on this design, from other things that trended over
the same decades: the Volcker monetary-policy regime change, the anchoring of
inflation expectations, changes in labor-market institutions and unionization,
globalization and import competition, changes in the frequency of price adjustment,
or a simple time trend / structural break. The most we can say is that \emph{the
low-frequency decline in the firm count and the decline in inflation's sensitivity to
the unemployment gap occurred together, with a sign consistent with the HSA theory}.
We therefore avoid ``competition causes flattening'' and ``we identify the
mechanism.''

\section{Conclusion}
Over 1982--2012 in the United States, the low-frequency decline in the number of
listed firms and the decline in inflation's sensitivity to the unemployment gap moved
together. The sign of the relationship---$\delta>0$, so that a falling firm count
implies a flatter slope---is consistent with the HSA prediction that rising market
concentration flattens the Phillips curve, and the implied $\kt$ falls from about
$0.12$ to near zero. Unifying inflation on core CPI produces the same sign and the
same flattening for the HP output gap. The start-to-end response difference for a
one-point unemployment gap is about 0.125 percentage point, an economically
nontrivial magnitude. A decomposition places the apparent
cross-measure disagreement in the 2008--09 episode, consistent with an important role
for energy prices. The relationship is identified mainly from samples that include
the 1980s; on post-1988 samples the slope has little variation left to explain and
the signal fades. The latent firm-count trend is nearly perfectly collinear with
linear time and assigns a much larger decline than the observed count; aggregate
markup movements correspond to about one quarter of the estimated flattening in a
simple calibration. Because of these facts, and because the empirical specification departs from the theory
in the ways set out in \autoref{sec:distance}, we read the finding as
\emph{theory-consistent time-series evidence, not a causal effect}: an early,
honest empirical operationalization of the HSA concentration--Phillips-curve
mechanism, whose value lies as much in showing where the evidence comes from as in
the headline sign.

\section{Status and open items}
\label{sec:status}
\paragraph{Confirmed.} Positive $\delta$ and a flattening $\kt$ for the unemployment
gap (priors/frequency/constraints/sub-samples); its economic magnitude; differences
among the observed count, source BN trend, and posterior latent trend, together with
period-specific identification shares; the core-CPI sign alignment for
output gaps (Bayesian and OLS/FWL); the concentration of the headline--core
disagreement in 2008--09; a well-identified $N$ decomposition (RMSE $\le 0.01$) and a
verified sampler (including missing-N annual-Q4).

\paragraph{Open (empirical/econometric).}
\begin{itemize}
\item \textbf{Causal identification.} A linear-time-trend horse race cannot distinguish
the two interactions. Add quadratic trends, structural-break dummies, and placebo
interactions before making any claim beyond association.
\item \textbf{Headline--core expectations mismatch.} Use core-matched expected
inflation, or a common trend-expectations series, or add explicit energy/import-price
controls, instead of treating the wedge as absorbed by the constant and $\alpha$.
\item \textbf{Predictive evaluation.} Replace the in-sample plug-in fit score with
LOO/WAIC or rolling out-of-sample forecasts.
\item \textbf{Period-output regeneration.} The period pipeline now records
\code{start}, \code{end}, and \code{n\_obs} after applying the same complete-case
filter used by the sampler, and applies the 40-observation minimum to that sample.
Existing period tables created before this correction must be regenerated; under
the corrected rule the 16-quarter post-2008 window is skipped.
\item \textbf{Bayes-factor robustness.} Report one-sided ($\delta>0$) assessments and
KDE-bandwidth sensitivity alongside the two-sided SDDR.
\end{itemize}

\paragraph{Open (computational).}
\begin{itemize}
\item \textbf{Chib marginal likelihood invalid for annual-Q4 blocks} (missing-N not
handled; $\log$-ML diverges). Fit scores and Savage--Dickey ratios are unaffected.
Fix: masked $N^{\mathrm{obs}}$ with \code{isfinite} branches in the Chib Kalman loops,
or blank the ML column for those blocks.
\item \textbf{\code{hsa\_full} weak convergence} (weak/tight/annual-Q4,
$\hat R\in[1.1,1.3]$): lengthen chains or drop those cells from primary inference.
\item \textbf{$\gamma$ unidentified} ($\hat N\!\cdot\!\bar N$ about $98\%$ collinear
with $\hat N$): treat \code{hsa\_const\_theta} as the default HSA representation.
\item Add the core-CPI specifications to the model-comparison pipeline so
\autoref{tab:mc} covers HP/BN core.
\end{itemize}

\paragraph{Key output files.} Coefficient posteriors
\code{results/tables/coefficient\_means.csv}; slope comparison
\code{results/tables/kappa\_model\_comparison.csv}; model comparison
\code{results/model\_comparison/}; figures \code{results/figures/}.

\appendix
\section{Appendix: pass-through from producer prices to core CPI}
\label{app:passthrough}
As an aggregate diagnostic of the pass-through decline suggested by HSA, we regress
core-CPI inflation on PPI inflation in rolling windows. ``Current'' reports the
current-PPI coefficient; ``distributed lag'' reports the sum from current PPI through
four quarterly lags. For 10-, 15-, and 20-year windows, the end-of-sample coefficient
is below its first available estimate.

\begin{table}[H]
\centering
\caption{Rolling OLS diagnostic of PPI-to-core-CPI pass-through. The correlation is
between each rolling estimate and the main specification's posterior mean $\bar N_t$.}
\label{tab:passthrough}
\begin{tabular}{llrrr}
\toprule
Window & PPI specification & First estimate & 2012Q4 & $\mathrm{corr}(\hat\rho_t,\bar N_t)$ \\
\midrule
10 years & Current & 0.091 & 0.007 & 0.641 \\
10 years & Current plus four lags, sum & 0.278 & 0.084 & 0.641 \\
15 years & Current & 0.062 & $-0.008$ & 0.724 \\
15 years & Current plus four lags, sum & 0.261 & 0.027 & 0.727 \\
20 years & Current & 0.075 & $-0.015$ & 0.732 \\
20 years & Current plus four lags, sum & 0.280 & 0.001 & 0.843 \\
\bottomrule
\end{tabular}
\end{table}

\noindent The decline in pass-through and its positive correlation with the firm-count
trend are directionally consistent with HSA, but this regression does not measure
firm-level price responses to cost shocks. PPI also reflects energy and import prices
and sectoral composition, while both the rolling coefficient and $\bar N_t$ carry
low-frequency trends. This is auxiliary descriptive mechanism evidence, not an
independent identification test for the main $\delta$.

\end{document}
